\section*{Capítulo \ref{ch:g3:shapes} --- Formas e ângulos retos}

\begin{solution}{exo:g3:shapes:1}
\omterm{def:g3:shapes:rightangle}{Ângulos retos} típicos: os cantos do quadro, de uma folha de papel,
da porta. Um ângulo que não é reto: a abertura de uma tesoura, os
ponteiros do relógio à 1 hora. Cada um é conferido com o esquadro
(\cref{met:g3:shapes:check}) --- nunca só no olho.
\end{solution}

\begin{solution}{exo:g3:shapes:2}
Construção: um lado de $6$ cm, \omterm{def:g3:shapes:rightangle}{ângulos retos} nas duas pontas com o
esquadro, lados de $4$ cm e fechar a forma. Marcas: $4$ quadradinhos
nos cantos, um tracinho em cada lado de $6$ cm e dois tracinhos em
cada lado de $4$ cm.
\end{solution}

\begin{solution}{exo:g3:shapes:3}
Um quadrado tem $4$ \omterm{def:g3:shapes:rightangle}{ângulos retos} e $4$ lados iguais --- todos
marcados.
\end{solution}

\begin{solution}{exo:g3:shapes:4}
Construções na malha; a malha garante os \omterm{def:g3:shapes:rightangle}{ângulos retos}, e contar
quadradinhos garante os comprimentos.
\end{solution}

\begin{solution}{exo:g3:shapes:5}
``Todo quadrado é um retângulo'': \emph{verdadeiro} --- ele tem $4$
\omterm{def:g3:shapes:rightangle}{ângulos retos}, que é tudo o que um retângulo exige.

``Todo retângulo é um quadrado'': \emph{falso} --- um retângulo
$6 \times 4$ tem lados desiguais.

``Um triângulo pode ter dois \omterm{def:g3:shapes:rightangle}{ângulos retos}'': \emph{falso} --- os
dois lados traçados a partir dos dois \omterm{def:g3:shapes:rightangle}{ângulos retos} seriam
perpendiculares à mesma base e, portanto, nunca se encontrariam para
fechar o triângulo.
\end{solution}

\begin{solution}{exo:g3:shapes:6}
$OA = 4$ cm (o \omterm{def:g3:shapes:shapes}{raio}). $OB = 4$ cm também: \emph{todo} ponto do
círculo está à mesma distância do centro.
\end{solution}

\begin{solution}{exo:g3:shapes:7}
O retângulo fica dividido em três triângulos: dois triângulos
retângulos ($AMD$ e $MBC$, com o \omterm{def:g3:shapes:rightangle}{ângulo reto} em $A$ e em $B$) e o
triângulo do meio $MCD$.
\end{solution}

\begin{solution}{exo:g3:shapes:8}
Com pelo menos um \omterm{def:g3:shapes:rightangle}{ângulo reto}: o quadrado, o retângulo, o triângulo
retângulo e a letra L (os cantos internos e externos dela). O
círculo não tem canto nenhum.
\end{solution}

\begin{solution}{exo:g3:shapes:9}
Cada lado percorre ``$2$ para a direita, $1$ para cima'' (ou as
versões giradas ``$1$ para a direita, $2$ para baixo'' etc.): cada
lado é a diagonal de um bloco $2 \times 1$ de quadradinhos, então os
quatro têm o mesmo comprimento.
\end{solution}

\begin{solution}{exo:g3:shapes:10}
Um programa correto: ``1. Trace um retângulo $ABCD$ com $AB = 3$ cm
(embaixo) e $BC = 2$ cm. 2. Sobre o lado $[DC]$, trace o quadrado
$DCEF$ de lado $3$ cm, do lado de fora do retângulo.'' Qualquer
programa claro, numerado e completo está correto --- teste-o
seguindo-o ao pé da letra.
\end{solution}

\begin{solution}{exo:g3:shapes:11}
Quadrados $1 \times 1$: $9$. Quadrados $2 \times 2$: $4$.
$3 \times 3$: $1$. Total: $9 + 4 + 1 = 14$ quadrados.
\end{solution}
